Copper production emissions are influenced by several interacting geological, technological, and energy-related factors. The complexity of the copper production chain implies that greenhouse gas emissions cannot be attributed to a single source alone. Instead, emissions result from the combined effects of production growth, ore quality, energy consumption, process characteristics, and the carbon intensity of the energy system. Understanding these different drivers is therefore essential for identifying potential technological and policy levers capable of reducing the environmental impact of copper mining.


\subsection{Production Growth}

One of the most direct drivers of CO$_2$ emissions is the total quantity of copper produced. As global demand for copper increases due to electrification and renewable energy deployment, mining operations must extract and process larger quantities of ore. The International Energy Agency projects that copper demand from clean energy technologies alone could more than double by 2040 under global net-zero scenarios \cite{IEA2021}. Higher production volumes generally require additional energy consumption across all stages of the production chain, including extraction, transportation, concentration, smelting, and refining processes. Consequently, production growth tends to increase both direct and indirect greenhouse gas emissions.

\subsection{Ore Grade Decline}

Ore grade represents the concentration of copper contained in extracted rock and constitutes one of the main structural constraints affecting mining sustainability. As ore grades decline, larger quantities of material must be extracted, transported, crushed, and processed in order to produce the same quantity of refined copper. Several studies report a long-term decline in average copper ore grades worldwide, with many operations now exploiting ores below 1\% copper concentration \cite{Northey2014,Calvo2016}. In Chile, average ore grades in concentrator plants declined from approximately 0.87\% in 2015 to around 0.66\% in 2024 \cite{Cochilco2024}. This increases energy requirements throughout the production chain and contributes directly to higher greenhouse gas emissions.

Unlike technological variables, ore grade is largely determined by geological conditions and resource depletion. Consequently, declining ore grades create a long-term structural pressure on mining energy consumption that may partially offset efficiency improvements achieved through technological progress.


\subsection{Energy Intensity of Mining Processes}

Copper production is highly energy-intensive due to the scale and complexity of mining and mineral processing operations. Different stages of the production chain require different forms and quantities of energy. Mining and transportation activities mainly rely on diesel-powered machinery, while crushing, grinding, concentration, electro-winning, and refining processes consume large amounts of electricity.

Existing studies show that concentrator plants are among the most energy-intensive stages of copper production. Cochilco reports electricity consumption levels exceeding 10,000 MJ per tonne of concentrate produced in certain concentration processes \cite{Cochilco2024}. At the national level, total energy consumption associated with Chilean copper mining increased by more than 15\% between 2015 and 2024 despite relatively stable production volumes. Energy intensity therefore constitutes a major determinant of CO$_2$ emissions.


\subsection{Electrification and Electricity Carbon Intensity}
\label{sec:elec and CI}
The distinction between direct and indirect emissions is particularly important in copper production. Direct emissions mainly originate from fuel combustion in extraction and transportation equipment, while indirect emissions are associated with electricity generation used throughout mineral processing operations.

Indirect emissions historically represented the largest share of emissions in Chilean copper mining due to the high electricity requirements of concentration and refining processes \cite{CochilcoEmissions2024}. However, the rapid decarbonization of Chile’s electricity system has significantly reduced the carbon intensity of electricity generation over recent years. Renewable sources represented nearly 70\% of Chile’s electricity generation in 2024 \cite{EnergiaEstrategica2025}. As a result, indirect emissions associated with electricity use in copper mining have declined substantially over the last decade.

Electrification of mining equipment is increasingly considered a major decarbonization pathway for the mining industry. However, the environmental benefits of electrification depend strongly on the carbon intensity of the electricity system. Consequently, future emissions reductions will depend both on technological changes within mining operations and on the continued decarbonization of electricity generation.


\subsection{Process Heterogeneity Along the Production Chain}
\label{sec:heterogeneity}
The different stages of copper production contribute unevenly to total energy consumption and greenhouse gas emissions. Certain operations, such as grinding, concentration, and smelting, are particularly energy-intensive, while others contribute comparatively less to total emissions. Moreover, the relative importance of fuel and electricity differs significantly between processes.

For example, mining and transportation stages rely predominantly on fuel combustion, whereas concentration and electro-winning processes are largely electricity-intensive \cite{Cochilco2024}. Understanding the heterogeneity of emissions across the production chain is therefore essential for identifying the most effective decarbonization strategies. Reducing emissions in processes with high energy consumption may have a significantly larger impact on total emissions than improving smaller contributors. Consequently, analysing the contribution of each process separately provides a more precise understanding of the technological and energy-related drivers influencing copper production emissions.